
\bibitem{mockus2002}
Audris Mockus, Roy T. Fielding, and James D. Herbsleb. 2002. Two case studies of open source software development: Apache and Mozilla. \emph{ACM Transactions on Software Engineering and Methodology} 11, 3, 309--346.

\bibitem{fielding1999}
Roy T. Fielding. 1999. Shared leadership in the Apache project. \emph{Communications of the ACM} 42, 4, 42--43.

\bibitem{crowston2005}
Kevin Crowston and James Howison. 2005. The social structure of free and open source software development. \emph{First Monday} 10, 2.

\bibitem{stewart2005}
Daniel Stewart. 2005. Social status in an open-source community. \emph{American Sociological Review} 70, 5, 823--842.

\bibitem{vonkrogh2003}
Georg von Krogh, Sebastian Spaeth, and Karim R. Lakhani. 2003. Community, joining, and specialization in open source software innovation: A case study. \emph{Research Policy} 32, 7, 1217--1241.

\bibitem{weber2004}
Steven Weber. 2004. \emph{The Success of Open Source}. Harvard University Press.

\bibitem{shah2006}
Sonali K. Shah. 2006. Motivation, governance, and the viability of hybrid forms in open source software development. \emph{Management Science} 52, 7, 1000--1014.

\bibitem{markus2007}
M. Lynne Markus. 2007. The governance of free/open source software projects: Monolithic, multidimensional, or configurational? \emph{Journal of Management \& Governance} 11, 2, 151--163.

\bibitem{dahlander2011}
Linus Dahlander and Siobh\'an O'Mahony. 2011. Progressing to the center: Coordinating project work. \emph{Organization Science} 22, 4, 961--979.

\bibitem{fleming2007}
Lee Fleming and David M. Waguespack. 2007. Brokerage, boundary spanning, and leadership in open innovation communities. \emph{Organization Science} 18, 2, 165--180.

\bibitem{shaikh2017}
Maha Shaikh and Ola Henfridsson. 2017. Governing open source software through coordination processes. \emph{Information and Organization} 27, 2, 116--135.

\bibitem{lindberg2016}
Aron Lindberg, Nicholas Berente, James Gaskin, and Kalle Lyytinen. 2016. Coordinating interdependencies in online communities: A study of an open source software project. \emph{Information Systems Research} 27, 4, 751--772.

\bibitem{howison2014}
James Howison and Kevin Crowston. 2014. Collaboration through open superposition: A theory of the open source way. \emph{MIS Quarterly} 38, 1, 29--50.

\bibitem{tsay2014}
Jason Tsay, Laura Dabbish, and James Herbsleb. 2014. Influence of social and technical factors for evaluating contribution in GitHub. In \emph{Proceedings of the 36th International Conference on Software Engineering (ICSE 2014)}, 356--366.

\bibitem{tsay2014b}
Jason Tsay, Laura Dabbish, and James Herbsleb. 2014. Let's talk about it: Evaluating contributions through discussion in GitHub. In \emph{Proceedings of the 22nd ACM SIGSOFT International Symposium on Foundations of Software Engineering (FSE 2014)}, 144--154.

\bibitem{marlow2013}
Jennifer Marlow, Laura Dabbish, and Jim Herbsleb. 2013. Impression formation in online peer production: Activity traces and personal profiles in GitHub. In \emph{Proceedings of the 2013 Conference on Computer Supported Cooperative Work (CSCW 2013)}, 117--128.

\bibitem{alami2020}
Adam Alami, Marisa Leavitt Cohn, and Andrzej W\k{a}sowski. 2020. How do FOSS communities decide to accept pull requests? In \emph{Proceedings of the 24th International Conference on Evaluation and Assessment in Software Engineering (EASE 2020)}, 220--229.

\bibitem{steinmacher2015}
Igor Steinmacher, Tayana Conte, Marco Aur\'elio Gerosa, and David Redmiles. 2015. Social barriers faced by newcomers placing their first contribution in open source software projects. In \emph{Proceedings of the 18th ACM Conference on Computer Supported Cooperative Work and Social Computing (CSCW 2015)}, 1379--1392.

\bibitem{dahlander2006}
Linus Dahlander and Martin W. Wallin. 2006. A man on the inside: Unlocking communities as complementary assets. \emph{Research Policy} 35, 8, 1243--1259.

\bibitem{schaarschmidt2015}
Mario Schaarschmidt, Gianfranco Walsh, and Harald F. O. von Kortzfleisch. 2015. How do firms influence open source software communities? A framework and empirical analysis of different governance modes. \emph{Information and Organization} 25, 2, 99--114.

\bibitem{riehle2014}
Dirk Riehle, Philipp Riemer, Carsten Kolassa, and Michael Schmidt. 2014. Paid vs.\ volunteer work in open source. In \emph{Proceedings of the 47th Hawaii International Conference on System Sciences (HICSS 2014)}, 3286--3295.

\bibitem{zhang2021}
Yuxia Zhang, Minghui Zhou, Audris Mockus, and Zhi Jin. 2021. Companies' participation in OSS development: An empirical study of OpenStack. \emph{IEEE Transactions on Software Engineering} 47, 10, 2242--2259.

\bibitem{zhang2022}
Yuxia Zhang, Klaas-Jan Stol, Hui Liu, and Minghui Zhou. 2022. Corporate dominance in open source ecosystems: A case study of OpenStack. In \emph{Proceedings of the 30th ACM Joint European Software Engineering Conference and Symposium on the Foundations of Software Engineering (ESEC/FSE 2022)}, 1048--1060.

\bibitem{butler2018}
Simon Butler, Jonas Gamalielsson, Bj\"orn Lundell, Christoffer Brax, Anders Mattsson, Tomas Gustavsson, Jonas Feist, Bengt Kvarnstr\"om, and Erik L\"onroth. 2018. An investigation of work practices used by companies making contributions to established OSS projects. In \emph{Proceedings of the 40th International Conference on Software Engineering: Software Engineering in Practice (ICSE-SEIP 2018)}, 201--210.

\bibitem{elliott2003}
Margaret S. Elliott and Walt Scacchi. 2003. Free software developers as an occupational community: Resolving conflicts and fostering collaboration. In \emph{Proceedings of the 2003 International ACM SIGGROUP Conference on Supporting Group Work (GROUP 2003)}, 21--30.

\bibitem{jensen2004}
Chris Jensen and Walt Scacchi. 2004. Collaboration, leadership, control, and conflict negotiation in the NetBeans.org community. In \emph{Proceedings of the 4th Workshop on Open Source Software Engineering (ICSE 2004)}, 48--52.

\bibitem{filippova2016}
Anna Filippova and Hichang Cho. 2016. The effects and antecedents of conflict in free and open source software development. In \emph{Proceedings of the 19th ACM Conference on Computer-Supported Cooperative Work and Social Computing (CSCW 2016)}, 705--716.

\bibitem{miller2022}
Courtney Miller, Sophie Cohen, Daniel Klug, Bogdan Vasilescu, and Christian K\"astner. 2022. ``Did you miss my comment or what?'' Understanding toxicity in open source discussions. In \emph{Proceedings of the 44th International Conference on Software Engineering (ICSE 2022)}, 710--722.

\bibitem{barcellini2008a}
Flore Barcellini, Fran\c{c}oise D\'etienne, and Jean-Marie Burkhardt. 2008. User and developer mediation in an open source software community: Boundary spanning through cross participation in online discussions. \emph{International Journal of Human-Computer Studies} 66, 7, 558--570.

\bibitem{barcellini2008b}
Flore Barcellini, Fran\c{c}oise D\'etienne, Jean-Marie Burkhardt, and Warren Sack. 2008. A socio-cognitive analysis of online design discussions in an open source software community. \emph{Interacting with Computers} 20, 1, 141--165.

\bibitem{sack2006}
Warren Sack, Fran\c{c}oise D\'etienne, Nicolas Ducheneaut, Jean-Marie Burkhardt, Dilan Mahendran, and Flore Barcellini. 2006. A methodological framework for socio-cognitive analyses of collaborative design of open source software. \emph{Computer Supported Cooperative Work} 15, 2--3, 229--250.

\bibitem{keertipati2016}
Smitha Keertipati, Sherlock A. Licorish, and Bastin Tony Roy Savarimuthu. 2016. Exploring decision-making processes in Python. In \emph{Proceedings of the 20th International Conference on Evaluation and Assessment in Software Engineering (EASE 2016)}, Article 43.

\bibitem{sharma2017}
Pankajeshwara Sharma, Bastin Tony Roy Savarimuthu, Nigel Stanger, Sherlock A. Licorish, and Austen Rainer. 2017. Investigating developers' email discussions during decision-making in Python language evolution. In \emph{Proceedings of the 21st International Conference on Evaluation and Assessment in Software Engineering (EASE 2017)}, 286--291.
